\section*{РОЗДІЛ 2. ПРОБЛЕМНЕ ПОЛЕ: КРИЗА РЕАКТИВНОЇ МОДЕЛІ БІБЛІОМЕТРИЧНОЇ ЗВІТНОСТІ}

\subsection*{2.1. Системні проблеми адміністрування науки в ЗВО}

Традиційна система адміністрування наукової діяльності у закладах вищої освіти формувалася в епоху, коли обсяг наукової продукції був на порядки меншим, а вимоги до звітності --- значно простішими. Однак за умов експоненційного зростання публікацій, різноманіття типів результатів (статті, препринти, дані, патенти) та посилення зовнішніх вимог до прозорості система поступово перетворюється на \textit{конвеєр регламентованої звітності}, який споживає ресурси, але слабо підтримує управлінські рішення.

Таким чином, криза реактивної моделі полягає не лише у «великому навантаженні на персонал», а у втраті керованості: дані збираються постфактум, з різних джерел і різної якості, а отже адміністративні процедури забезпечують здебільшого \textbf{compliance}, а не \textbf{performance management} та стратегічний розвиток. Далі розглянуто ключові механізми цієї кризи.

\subsubsection*{2.1.1. Фрагментація даних та відсутність «єдиного джерела правди»}

Фундаментальною передумовою кризового стану є фрагментація наукометричних даних: інформація про одного й того самого дослідника та його результати існує як набір частково перекривних «проєкцій» у різних системах. На практиці це породжує щонайменше чотири розриви.

\begin{enumerate}[leftmargin=*,itemsep=0.2em]
\item \textbf{Розрив ідентичності:} відмінності в іменуванні авторів, афіліаціях, наявності/відсутності ідентифікаторів.
\item \textbf{Розрив об'єкта:} та сама публікація репрезентується різними наборами метаданих (назва, журнал, DOI, рік, тип документа).
\item \textbf{Розрив доступу й автоматизації:} неоднакові можливості машинного доступу (наявність API, обмеження ліцензій, відсутність офіційного API).
\item \textbf{Розрив охоплення:} міжнародні індекси, відкриті графи знань та локальні реєстри фіксують різні сегменти наукового доробку.
\end{enumerate}

Ці розриви підсилюються тим, що інституції одночасно покладаються на кілька джерел: \textbf{ORCID} як реєстр ідентичностей і афіліацій (19+ млн профілів станом на 2024 рік \cite{ORCID2024APIDocs}); \textbf{OpenAlex} як відкрита база метаданих (250+ млн праць, 100+ млн авторів \cite{Priem2022OpenAlex,OpenAlex2024Docs}); \textbf{Scopus} та \textbf{Web of Science} як комерційні індекси з різним покриттям і ліцензійними обмеженнями; \textbf{Google Scholar} як найбільш «широке», але слабко автоматизоване джерело; а також \textbf{локальні бази ЗВО} (часто Excel або застарілі БД), де зберігаються дані, необхідні для внутрішньої звітності.

За відсутності інтеграції цієї екосистеми інституція не може сформувати \textbf{«Golden Record»} (єдиний авторитетний запис) для кожного дослідника й кожного результату. У результаті звітність системно стикається з трьома типами дефектів.

\textbf{Дублювання записів.} Одна публікація потрапляє в звіт кілька разів (наприклад, з ORCID, Scopus і локальної бази) з незначними відмінностями в метаданих. За нашими оцінками (вибірковий аналіз типових річних звітів кафедр), \textbf{15--25\% записів є дублікатами}.

\textbf{Неузгодженість метаданих.} Різні бази фіксують різні значення цитувань, ключових слів чи афіліацій, тому питання «який запис правильний» часто вирішується вручну й постфактум.

\textbf{Неповнота охоплення.} Жодне джерело не дає повної картини: частина публікацій у вітчизняних виданнях може бути відсутня у Scopus/WoS; препринти та тези не завжди системно відображені; монографії й локальні видання часто залишаються поза полем міжнародних індексів.

\subsubsection*{2.1.2. Трудомісткість та ресурсоємність звітності}

Фрагментація даних безпосередньо трансформується у витрати: що більше джерел та форматів, то більше ручних узгоджень і перевірок. На основі аналізу практик українських ЗВО можна оцінити типові часові затрати на підготовку регламентованих документів (табл. 2.1), які фактично є «податком на відсутність інтегрованих даних».

\begin{longtable}{@{}p{0.30\textwidth}p{0.24\textwidth}p{0.40\textwidth}@{}}
\toprule
\textbf{Тип звіту} & \textbf{Традиційний час підготовки} & \textbf{Основні операції} \\
\midrule
Річний звіт кафедри & 8--12 годин & Збір даних з 3--5 джерел, верифікація, форматування \\
Атестація науково-педагогічного працівника & 3--5 годин & Компіляція списку публікацій, підрахунок показників \\
Конкурсна документація & 5--8 годин & Оформлення за специфічними вимогами конкурсу \\
Звіт за грантовим проєктом & 4--6 годин & Агрегація показників учасників, звірка з планом \\
\bottomrule
\end{longtable}

Як правило, процес підготовки звіту проходить повторюваний «ланцюжок» етапів: збір фрагментів інформації, їх звірка й дедуплікація, нормалізація бібліографічного опису, обчислення показників, фінальний контроль якості. Левову частку часу поглинають такі ручні операції.

\begin{enumerate}[leftmargin=*,itemsep=0.2em]
\item \textit{Копіювання та вставка} (copy-paste) даних з веб-інтерфейсів наукометричних баз у текстові документи чи таблиці.
\item \textit{Верифікація та дедуплікація} --- ручна перевірка, чи не внесено одну публікацію двічі, звірка написання імен авторів, журналів.
\item \textit{Форматування} --- приведення бібліографічних описів до вимог ДСТУ, МОН, конкретного конкурсу.
\item \textit{Підрахунок показників} --- ручний розрахунок H-індексів, сум цитувань, розподілу за роками та квартилями.
\item \textit{Контроль якості} --- фінальна вичитка на предмет помилок, пропусків, невідповідностей.
\end{enumerate}

Критично, ці операції не лише споживають час висококваліфікованих працівників, який міг би бути спрямований на аналітику та розвиток, а й створюють \textit{системний канал помилок}: при обробці сотень записів з різною структурою неминуче зростають пропуски, арифметичні неточності та порушення форматування.

\subsubsection*{2.1.3. Реактивність та низька аналітична цінність}

Якщо фрагментація пояснює, \textit{чому} дані важко зібрати, то реактивність описує, \textit{як} ці дані вбудовані в цикл управління. У традиційній моделі звіти формуються post factum у відповідь на зовнішній запит (наказ МОН, вимога ректорату, дедлайн конкурсу), тобто звітність працює як механізм фіксації минулого, а не як інструмент підтримки рішень у режимі \textit{near real-time}. У такій конфігурації «управління» підміняється інвентаризацією: спочатку організація витрачає ресурси на зведення, а вже потім (за залишковим принципом) на інтерпретацію.

Наслідки цієї реактивності проявляються не точково, а системно.

\textbf{Оперативне реагування на тренди.} Коли основний індикатор продуктивності формується раз на рік, керівництво кафедри або ЗВО дізнається про зміну публікаційної активності із запізненням до 12 місяців. Вікно для коригувальних дій звужується, а ранні сигнали (наприклад, зменшення кількості прийнятих до друку рукописів) залишаються поза управлінським полем.

\textbf{Прогнозування та планування.} Планування потребує не лише історичних підсумків, а й динаміки: потоку оновлень про результати, грантові заявки, співавторства та афіліації. За відсутності «Golden Record» і стабільної верифікації метаданих прогнози перетворюються на експертні припущення, а не на відтворювані розрахунки.

\textbf{Бенчмаркінг.} Регулярне порівняння з іншими підрозділами можливе лише тоді, коли показники є зіставними за методикою та джерелами. У реактивній моделі кожне порівняння запускає новий цикл збору й ручної нормалізації, тому бенчмаркінг стає разовою активністю, а не практикою.

\textbf{Виявлення прихованих патернів.} Коли значна частина часу витрачається на очищення й узгодження, аналітика редукується до описових підсумків. Відповіді на питання «чому так сталося» (структурні причини) та «що робити далі» (policy/management interventions) залишаються за межами регламентованої звітності.

Отже, традиційна система виконує функцію \textbf{compliance}, а не \textbf{strategy}: вона гарантує формальну відповідність вимогам, але не забезпечує інструментів управління на основі даних. Далі показано, чому саме український контекст робить ці обмеження особливо виразними.

\subsection*{2.2. Контекст українських ЗВО: специфіка та можливості}

Описані вище механізми є загальними для багатьох освітньо-наукових систем, однак в українських ЗВО вони посилюються одночасною дією трьох факторів: \textit{(1) регуляторного тиску, (2) євроінтеграційних зобов'язань, (3) ресурсних обмежень}. Разом вони створюють ситуацію, коли інституції змушені нарощувати формальну звітність швидше, ніж здатні модернізувати дані та інструменти їх обробки.

\subsubsection*{2.2.1. Регуляторне середовище}

Українські ЗВО функціонують у специфічному регуляторному середовищі, що визначається вимогами Міністерства освіти і науки України. Ці вимоги задають високий рівень деталізації звітних форм і, відповідно, роблять чутливими будь-які помилки в метаданих.

Ключові документи:
\begin{itemize}[leftmargin=*,itemsep=0.2em]
\item \textbf{Порядок формування Переліку наукових фахових видань України} \cite{MONUkraine2023ProfessionalJournalsOrder} визначає категорії журналів (А, Б) та вимоги до індексації в міжнародних базах;
\item вимоги до звітності наукових підрозділів включають детальну інформацію про публікації, гранти, захисти дисертацій, патенти;
\item критерії ліцензування та акредитації освітніх програм включають наукометричні показники викладачів.
\end{itemize}

У такій системі навіть незначні розбіжності між джерелами (назва журналу, афіліація, тип документа) можуть мати адміністративні наслідки, а отже стимулюють додаткові раунди ручної перевірки.

\subsubsection*{2.2.2. Євроінтеграційний контекст}

Євроінтеграційний вектор, з одного боку, відкриває можливості для стандартів і інтероперабельності, а з іншого --- підвищує планку вимог до якості даних. Інтеграція України до Європейського дослідницького простору (European Research Area, ERA) передбачає відповідність стандартам OpenAIRE та EOSC \cite{OpenAIRE2023GuidelinesCRIS}. Це включає:

\begin{itemize}[leftmargin=*,itemsep=0.2em]
\item підтримку стандарту \textbf{CERIF} (Common European Research Information Format) для обміну дослідницькою інформацією \cite{euroCRIS2024CERIF};
\item інтеграцію інституційних репозитаріїв з European Open Science Cloud;
\item забезпечення відкритого доступу до результатів досліджень, що фінансуються з публічних коштів.
\end{itemize}

Однак більшість українських ЗВО досі не мають повноцінних CRIS-систем, що відповідають цим стандартам, покладаючись на комбінацію застарілих локальних баз та ручної обробки. У результаті розрив між «стандартами обміну» і «практиками ведення даних» зберігається.

\subsubsection*{2.2.3. Ресурсні обмеження та можливості}

Ресурсний вимір визначає, які саме траєкторії модернізації є реалістичними для інституцій.

\textbf{Обмеження:}
\begin{itemize}[leftmargin=*,itemsep=0.2em]
\item недостатнє фінансування для придбання комерційних CRIS-платформ (Pure, Converis коштують десятки тисяч євро на рік);
\item дефіцит кваліфікованих ІТ-кадрів для розробки та підтримки власних рішень;
\item низький рівень цифрової грамотності частини адміністративного персоналу.
\end{itemize}

\textbf{Можливості:}
\begin{itemize}[leftmargin=*,itemsep=0.2em]
\item партнерство Міністерства цифрової трансформації України з University College London у сфері AI governance створює основу для трансферу знань;
\item інвестиції World Bank та інших донорів у цифрову інфраструктуру освіти;
\item доступність відкритих API (OpenAlex, ORCID) та open-source рішень (DSpace-CRIS);
\item поява нових фахових видань з фокусом на цифрових практиках публічного управління («Публічно-управлінські та цифрові практики», 2025).
\end{itemize}

Отже, українські ЗВО опиняються в конфігурації, де «дорогі» інституційні рішення часто недоступні, але відкриті дані та інструменти створюють простір для альтернативних архітектур (інтеграції через API, агентні підходи, поступова нормалізація даних).

\subsubsection*{2.2.4. Формулювання дослідницького gap}

Синтезуючи наведені аргументи, можна сформулювати \textbf{дослідницький розрив (research gap)}: попри наявність CRIS-підходів, локальних регламентів та окремих інтеграцій, на практиці вони не усувають ключовий ланцюжок проблем \textit{(фрагментація $\rightarrow$ ручна нормалізація $\rightarrow$ реактивність $\rightarrow$ низька аналітична цінність)} в умовах обмежених ресурсів українських ЗВО.

Звідси випливають конкретні наукові й прикладні питання: \textit{якою має бути архітектура та методи обробки даних, щоб зменшити частку ручної праці, підвищити відтворюваність показників і перевести звітність у режим проактивного моніторингу?}

Відповідно, існуючі наукові розвідки та практичні рішення не забезпечують:
\begin{enumerate}[leftmargin=*,itemsep=0.2em]
\item \textbf{Концептуальної моделі} переходу від реактивної до проактивної парадигми адміністрування науки з урахуванням можливостей GenAI;
\item \textbf{Механізмів автономної обробки} мультиджерельних наукометричних даних, де ШІ-агент самостійно визначає послідовність дій (на відміну від жорстко запрограмованих ETL-процесів);
\item \textbf{Алгоритмів інтелектуального злиття} даних з різних джерел із формуванням Golden Record та трекінгом provenance;
\item \textbf{Емпіричної валідації} ефективності GenAI-рішень у реальних умовах українських ЗВО з вимірюванням кількісних показників.
\end{enumerate}

Заповнення цього розриву є метою нашого дослідження і визначає логіку наступних розділів, у яких обґрунтовуються концептуальна модель та методичні рішення.